% ============================================================================
% Fichier  : partie2/P02_C07_vulnerabilites.tex
% Titre    : Analyser et Prioriser les Vulnérabilités
% Auteur   : MINKA MI NGUIDJOÏ Thierry Emmanuel
% Version  : 1.0 -- novembre 2026
% ============================================================================
% Sources : FIRST, CVSS v3.1 et EPSS ; CISA, catalogue Known Exploited
% Vulnerabilities ; programme CVE (MITRE) ; NIST SP 800-40 Rev. 4.
% Résultats de l'Opération LOGONE : cas fictif.
% ============================================================================

\chapter{Analyser et Prioriser les Vulnérabilités}
\label{chap:vulnerabilites}
\chaptermark{Analyser et prioriser}

\epigraph{%
    « Un rapport de trois cents failles\\
    est un rapport qu'on ne lira pas. »%
}{Principe du cours}

\niveauChapitre{Intermédiaire}{Chapitres~\ref{chap:methodes}
et~\ref{chap:balayage}.}

\begin{boxobjectifs}
\begin{itemize}
    \item Utiliser un scanner de vulnérabilités et en connaître les limites.
    \item Écarter les faux positifs par la vérification.
    \item Comprendre les identifiants CVE et les catalogues de failles
          exploitées.
    \item Prioriser par la gravité, l'exploitation connue, l'exposition et
          la valeur de l'actif.
    \item Présenter une liste courte, ordonnée et justifiée.
\end{itemize}
\end{boxobjectifs}

\bigskip

% ============================================================================
\section*{Question directrice : par quoi commencer, parmi 312 résultats~?}
% ============================================================================

Le scanner lancé le 26~octobre sur les réseaux de Garoua produit 312
résultats. M.~KOUNOU n'a ni le temps ni l'équipe de tout corriger. Mme~NGO
MBOCK doit lui remettre une liste courte~: ce qui est vrai, ce qui est
dangereux, et dans quel ordre.

\medskip
\textbf{Plan du chapitre.}
\begin{enumerate}
    \item Les scanners de vulnérabilités (section~\ref{sec:c07-scanners}).
    \item Vérifier~: faux positifs et faux négatifs (section~\ref{sec:c07-verifier}).
    \item CVE et catalogues (section~\ref{sec:c07-cve}).
    \item Prioriser (section~\ref{sec:c07-prioriser}).
    \item Présenter (section~\ref{sec:c07-presenter}).
\end{enumerate}

\separateur

% ============================================================================
\section{Les Scanners de Vulnérabilités}
\label{sec:c07-scanners}
% ============================================================================

Un \termecle{scanner de vulnérabilités} compare les services et les
versions observés à une base de failles connues, et tente parfois des
vérifications. Il peut travailler \emph{sans authentification} (vue
d'attaquant externe) ou \emph{avec authentification} (il lit les versions
installées~: bien plus précis).

\begin{tableaucours}{L{3.6cm}L{8.6cm}}{%
    \tableauHeader{Limite} &
    \tableauHeader{Conséquence}
}
Déduction par la version & Une version corrigée par un rétroportage peut
être signalée vulnérable \\
Bases incomplètes & Les failles de configuration et de logique sont mal vues \\
Absence de contexte & Le scanner ne sait pas ce qui est critique pour le
métier \\
Bruit & Des centaines de résultats de faible intérêt \\
\end{tableaucours}
\vspace{-0.8em}
\begin{center}{\small\itshape Tableau~7.1. Limites des scanners}\end{center}

% ============================================================================
\section{Vérifier : Faux Positifs et Faux Négatifs}
\label{sec:c07-verifier}
% ============================================================================

\begin{boxdef}
\termecle{Faux positif}~: faille signalée qui n'existe pas.
\termecle{Faux négatif}~: faille présente que l'outil n'a pas signalée. Le
premier fait perdre du temps~; le second donne une fausse assurance.
\end{boxdef}

\begin{boxscenario}{Opération LOGONE. De 312 à 41}
\textbf{Constatations.} Sur 312~résultats bruts~: 158 sont des
informations sans gravité~; 96 sont des doublons (même faille sur plusieurs
hôtes, regroupée)~; 17 sont des faux positifs (versions corrigées par
rétroportage). Restent 41~vulnérabilités distinctes et vérifiées, dont 3
figurent au catalogue des failles exploitées.

\medskip
\textbf{Analyse} (confiance élevée). Le travail du testeur n'est pas de
recopier le scanner, mais de transformer ses résultats en constats vrais et
ordonnés.
\end{boxscenario}

% ============================================================================
\section{CVE et Catalogues}
\label{sec:c07-cve}
% ============================================================================

\begin{tableaucours}{L{3.4cm}L{8.8cm}}{%
    \tableauHeader{Référence} &
    \tableauHeader{Rôle}
}
CVE & Identifiant public unique d'une vulnérabilité (programme CVE, MITRE)~\cite{mitrecve} \\
CVSS & Gravité technique, de 0 à 10 (FIRST)~\cite{first2019cvss31} \\
EPSS & Probabilité estimée d'exploitation dans les 30~jours (FIRST)~\cite{firstepss} \\
Catalogue KEV & Liste des vulnérabilités dont l'exploitation est avérée (CISA)~\cite{cisakev} \\
\end{tableaucours}
\vspace{-0.8em}
\begin{center}{\small\itshape Tableau~7.2. Les références utiles à la priorisation}\end{center}

% ============================================================================
\section{Prioriser}
\label{sec:c07-prioriser}
% ============================================================================

\begin{boxprocedure}
\textbf{Une règle de priorisation (convention du cours)}
\begin{enumerate}
    \item \textbf{Priorité~1}~: exploitation avérée (catalogue KEV) et actif
          exposé à Internet ou critique~; ou CVSS d'au moins 9 sur un actif
          exposé.
    \item \textbf{Priorité~2}~: CVSS d'au moins 7 sur un actif interne
          critique, ou probabilité d'exploitation élevée.
    \item \textbf{Priorité~3}~: le reste, traité dans le cycle normal des
          correctifs (cours SEC3266, Chapitre sur les systèmes).
\end{enumerate}
\end{boxprocedure}

\begin{boxalert}
\textbf{Violation typique.} Trier uniquement par le score CVSS décroissant.

\smallskip
\textbf{Règle~:} une faille de score 7,5 exploitée activement sur la
passerelle passe avant une faille de score 9,1 sur un poste isolé, sans
exploitation connue. On croise gravité, exploitation, exposition et valeur.
\end{boxalert}

% ============================================================================
\section{Présenter}
\label{sec:c07-presenter}
% ============================================================================

\begin{boxobjet}{T-06 · SMBv1 actif sur deux serveurs de Garoua}
\textbf{Description~:} le protocole SMBv1, obsolète, est accepté par
FS-02 et APP-01. \textbf{Preuve~:} négociation du protocole (relevé du
26~octobre). \textbf{Gravité~:} élevée (vulnérabilités connues de SMBv1,
exploitées par le passé à grande échelle). \textbf{Impact~:} exécution de
code à distance si un correctif manque. \textbf{Recommandation~:}
désactiver SMBv1, vérifier les clients qui en dépendent.
\textbf{Détection~:} journaliser les négociations SMBv1.
\end{boxobjet}

\begin{boxmetier}{RSSI}
\textbf{Demander peu, mais juste.} Le RSSI attend du rapport une liste de
priorité~1 courte et argumentée, qu'il peut faire corriger dans le mois, et
une annexe complète pour le cycle normal.
\end{boxmetier}

\separateur

\begin{boxresume}
\textbf{Ce qu'il faut retenir de ce chapitre~:}
\begin{itemize}
    \item \textbf{Un scanner} compare des versions à une base~; avec
          authentification, il est plus précis.
    \item \textbf{On vérifie}~: faux positifs écartés, faux négatifs
          recherchés à la main.
    \item \textbf{CVE, CVSS, EPSS, KEV}~: identifier, coter, estimer,
          savoir ce qui est exploité.
    \item \textbf{On priorise} en croisant gravité, exploitation, exposition
          et valeur.
    \item \textbf{On présente} une liste courte de priorité~1 et une annexe.
\end{itemize}
\end{boxresume}

\begin{boxinfo}
\textbf{Points de contrôle de l'élève.} Je sais~: dire les limites d'un
scanner~; distinguer faux positif et faux négatif~; expliquer CVE, CVSS,
EPSS et KEV~; appliquer une règle de priorisation~; rédiger un constat.
\end{boxinfo}

\bigskip

\begin{boxexo}[title={\ding{45}\quad Exercice~7.1. Les chiffres \niveaubase}]
Vérifiez que 312 résultats moins les informations, les doublons et les faux
positifs laissent 41~vulnérabilités.
\end{boxexo}

\begin{boxexo}[title={\ding{45}\quad Exercice~7.2. Faux positif \niveaubase}]
Pourquoi une version corrigée par rétroportage est-elle souvent signalée à
tort~?
\end{boxexo}

\begin{boxexo}[title={\ding{45}\quad Exercice~7.3. Les références \niveaubase}]
Associez à chaque question la bonne référence (CVE, CVSS, EPSS, KEV)~:
« est-elle exploitée~? »~; « quelle gravité~? »~; « quel nom~? »~; « quelle
probabilité d'être exploitée bientôt~? »
\end{boxexo}

\begin{boxexo}[title={\ding{45}\quad Exercice~7.4. Prioriser \niveaumoyen}]
Appliquez la règle du cours à~: CVSS 7,5 au catalogue KEV sur la
passerelle~; CVSS 9,1 sans exploitation sur un poste isolé~; CVSS 5,3 sur le
portail.
\end{boxexo}

\begin{boxexo}[title={\ding{45}\quad Exercice~7.5. Authentifié \niveaumoyen}]
Pourquoi un balayage authentifié réduit-il les faux positifs~? Quelles
précautions exige-t-il~?
\end{boxexo}

\begin{boxexo}[title={\ding{45}\quad Exercice~7.6. Faux négatifs \niveaumoyen}]
Citez trois types de failles qu'un scanner voit mal, et comment le testeur
les recherche.
\end{boxexo}

\begin{boxexo}[title={\ding{45}\quad Exercice~7.7. T-06 \niveaumoyen}]
Quels clients pourraient encore dépendre de SMBv1 chez LiZbiZ~? Comment le
vérifier avant de le désactiver~?
\end{boxexo}

\begin{boxexo}[title={\ding{45}\quad Exercice~7.8. Regrouper \niveaumoyen}]
Pourquoi regrouper une même faille présente sur 40~hôtes en un seul
constat~? Que faut-il indiquer pour ne rien perdre~?
\end{boxexo}

\begin{boxexo}[title={\ding{45}\quad Exercice~7.9. Une politique \niveauexpert}]
Rédigez la politique de priorisation des vulnérabilités de LiZbiZ, avec les
délais associés à chaque priorité.
\end{boxexo}

\begin{boxexo}[title={\ding{45}\quad Exercice~7.10. Mesurer \niveauexpert}]
Proposez trois indicateurs du processus de gestion des vulnérabilités et
leurs seuils.
\end{boxexo}

\bigskip

\begin{boxinfo}
\textbf{Pour aller plus loin.}
\begin{itemize}
    \item FIRST, CVSS~\cite{first2019cvss31} et EPSS~\cite{firstepss}~;
          CISA, catalogue KEV~\cite{cisakev}.
    \item NIST SP~800-40 Rev.~4~\cite{nist80040}.
\end{itemize}
\end{boxinfo}

\begin{boxinfo}
\textbf{Transition.} Une faille vérifiée se démontre par l'exploitation,
dans un cadre maîtrisé. La Partie~III commence par ses principes
(Chapitre~\ref{chap:exploitation}).
\end{boxinfo}
